\subsection{Convolution}
\label{jep104local033}

We will make extensive use of the convolution construction, defined as follows. Consider $\jepPubEuler{F},\jepPubEuler{G}$ in $D^{\mathrm{b}}_{G_{\jepPubScript{O}}}(\mathrm{Gr},\Bbbk)$, and the diagram
\[
 \mathrm{Gr} \times \mathrm{Gr} \xleftarrow{\textstyle p^\mathrm{Gr}} G_{\jepPubScript{K}} \times \mathrm{Gr} \xrightarrow{\textstyle q^\mathrm{Gr}} G_{\jepPubScript{K}} \times^{G_{\jepPubScript{O}}} \mathrm{Gr} \xrightarrow{\textstyle m^\mathrm{Gr}} \mathrm{Gr},
\]
where $p^\mathrm{Gr}$ and $q^\mathrm{Gr}$ are the quotient morphisms, and $m^\mathrm{Gr}$ is induced by the $G_{\jepPubScript{K}}$-action on $\mathrm{Gr}$. Consider the action of $G_{\jepPubScript{O}} \times G_{\jepPubScript{O}}$ on $G_{\jepPubScript{K}} \times \mathrm{Gr}$ defined by
\[
 (g_1,g_2) \cdot (h_1,h_2G_{\jepPubScript{O}}) = (g_1h_1(g_2)^{-1}, g_2 h_2G_{\jepPubScript{O}}).
\]
Then the functor $(q^\mathrm{Gr})^*$ induces an equivalence of categories
\[
 D^{\mathrm{b}}_{G_{\jepPubScript{O}}}(G_{\jepPubScript{K}} \times^{G_{\jepPubScript{O}}} \mathrm{Gr},\Bbbk) \xrightarrow{\sim} D^{\mathrm{b}}_{G_{\jepPubScript{O}} \times G_{\jepPubScript{O}}}(G_{\jepPubScript{K}} \times \mathrm{Gr},\Bbbk).
\]
Hence there exists a unique object $\jepPubEuler{F} \, \widetilde{\boxtimes} \, \jepPubEuler{G}$ such that
\[
 (q^\mathrm{Gr})^* \bigl( \jepPubEuler{F} \, \widetilde{\boxtimes} \, \jepPubEuler{G} \bigr) = (p^\mathrm{Gr})^* \bigl( \jepPubEuler{F} \jepDerivedExternalTensor_{\Bbbk} \jepPubEuler{G} \bigr).
\]
The convolution product of $\jepPubEuler{F}$ and $\jepPubEuler{G}$ is defined by
\begin{equation}
\label{jep104local008}
\jepPubEuler{F} \star^{G_{\jepPubScript{O}}} \jepPubEuler{G} := (m^\mathrm{Gr})_* \bigl( \jepPubEuler{F} \, \widetilde{\boxtimes} \, \jepPubEuler{G} \bigr).
\end{equation}
This construction endows the category $D^{\mathrm{b}}_{G_{\jepPubScript{O}}}(\mathrm{Gr},\Bbbk)$ with the structure of a monoidal category. A similar formula defines a right action of this monoidal category on $D^{\mathrm{b}}_K(\mathrm{Gr},\Bbbk)$, for any $K \subset G_{\jepPubScript{O}}$. (This action will again be denoted $\star^{G_{\jepPubScript{O}}}$.)

\begin{rmk}
 Note that if $\Bbbk$ is not a field, the convolution product considered above is \emph{not} the same as the one considered (when $\jepPubEuler{F}$ and $\jepPubEuler{G}$ are perverse sheaves) in~\cite{jep104ref36}: the product considered in~\cite{jep104ref36} is rather defined as ${}^{p}\jepPubEuler{H}^0(\jepPubEuler{F} \star^{G_{\jepPubScript{O}}} \jepPubEuler{G})$ in our notation.
\end{rmk}

\begin{lem}
\label{jep104local018}
 Assume that $\Bbbk$ is a field. If $\jepPubEuler{F}$ belongs to $\mathrm{Perv}(\mathrm{Gr},\Bbbk)$ and $\jepPubEuler{G}$ belongs to $\mathrm{Perv}_{G_{\jepPubScript{O}}}(\mathrm{Gr},\Bbbk)$, then $\jepPubEuler{F} \star^{G_{\jepPubScript{O}}} \jepPubEuler{G}$ belongs to $\mathrm{Perv}(\mathrm{Gr},\Bbbk)$.
\end{lem}

\begin{proof}
 This claim follows from the description of convolution in terms of nearby cycles obtained in~\cite[Prop.~6]{jep104ref27}. (In~\cite{jep104ref27}, only the case of characteristic-$0$ coefficients is treated. However the same proof applies in general, simply replacing~\cite[Prop.~1]{jep104ref27} by~\cite[Prop.~2.2]{jep104ref36}.)
\end{proof}

\begin{rmk}
The description of convolution in terms of nearby cycles as in~\cite{jep104ref27} works for general coefficients (if convolution is defined as in~\eqref{jep104local008}). The nearby cycles functor is $t$-exact in this generality, but this description also involves a (derived) external tensor product. If $\Bbbk$ is not a field this tensor product operation is not $t$-exact, which explains the failure of Lemma~\ref{jep104local018} in this setting.
\end{rmk}

A very similar construction as the one considered above, based on the diagram
\[
 \mathrm{Fl} \times \mathrm{Fl} \xleftarrow{\textstyle p^\mathrm{Fl}} G_{\jepPubScript{K}} \times \mathrm{Fl} \xrightarrow{\textstyle q^\mathrm{Fl}} G_{\jepPubScript{K}} \times^{I^-} \mathrm{Fl} \xrightarrow{\textstyle m^\mathrm{Fl}} \mathrm{Fl},
\]
provides a convolution product $\star^{I^-}$ on $D^{\mathrm{b}}_{I^-}(\mathrm{Fl},\Bbbk)$, which endows this category with the structure of a monoidal category, and defines a right action of this monoidal category on $D^{\mathrm{b}}_K(\mathrm{Fl},\Bbbk)$, for any $K \subset G_{\jepPubScript{O}}$.

Again the same formulas, using the diagram
\[
 \mathrm{Fl} \times \mathrm{Gr} \xleftarrow{\textstyle p^{\mathrm{Fl}}_{\mathrm{Gr}}} G_{\jepPubScript{K}} \times \mathrm{Gr} \xrightarrow{\textstyle q^{\mathrm{Fl}}_{\mathrm{Gr}}} G_{\jepPubScript{K}} \times^{I^-} \mathrm{Gr} \xrightarrow{\textstyle m^{\mathrm{Fl}}_{\mathrm{Gr}}} \mathrm{Gr},
\]
allows to define a bifunctor
\[
 D^{\mathrm{b}}_K(\mathrm{Fl},\Bbbk) \times D^{\mathrm{b}}_{I^-}(\mathrm{Gr},\Bbbk) \to D^{\mathrm{b}}_K(\mathrm{Gr},\Bbbk),
\]
which will once again be denoted $\star^{I^-}$.

The following lemma is standard; its proof is left to interested readers.

\begin{lem}
\label{jep104local017}
 For any subgroup $K\subset G_{\jepPubScript{O}}$, any $\jepPubEuler{F}$ in $D^{\mathrm{b}}_K(\mathrm{Fl},\Bbbk)$ and any $\jepPubEuler{G}$ in $D^{\mathrm{b}}_{G_{\jepPubScript{O}}}(\mathrm{Gr},\Bbbk)$, there exists a canonical isomorphism
 \[
  \jepPubEuler{F} \star^{I^-} \mathsf{For}^{G_{\jepPubScript{O}}}_{I^-}(\jepPubEuler{G}) \cong \pi_*(\jepPubEuler{F}) \star^{G_{\jepPubScript{O}}} \jepPubEuler{G}
 \]
 in $D^{\mathrm{b}}_K(\mathrm{Gr},\Bbbk)$.
\end{lem}

In the following lemma we consider the convolution bifunctor
\[
 (-) \star^{G_{\jepPubScript{O}}} (-) : D^{\mathrm{b}}_K(\mathrm{Gr},\Bbbk) \times D^{\mathrm{b}}_{G_{\jepPubScript{O}}}(\mathrm{Fl},\Bbbk) \to D^{\mathrm{b}}_K(\mathrm{Fl},\Bbbk)
\]
(constructed once again using formulas similar to those above). Its proof is easy, and

left to the reader.

\begin{lem}
\label{jep104local019}
 Let $\jepPubEuler{F}$ in $D^{\mathrm{b}}_K(\mathrm{Gr},\Bbbk)$ and $\jepPubEuler{G}$ in $D^{\mathrm{b}}_{I^-}(\mathrm{Fl},\Bbbk)$. Then there exists a canonical isomorphism
 \[
  \pi^*(\jepPubEuler{F}) \star^{I^-} \jepPubEuler{G} \cong \jepPubEuler{F} \star^{G_{\jepPubScript{O}}} {}^{*}\mathsf{Ind}_{I^-}^{G_{\jepPubScript{O}}}(\jepPubEuler{G})
 \]
 in $D^{\mathrm{b}}_K(\mathrm{Fl},\Bbbk)$.
\end{lem}

